\section{Conclusion}
\label{sec:conclusion}

We formulated ANN-to-SNN conversion for graph node classification and found that the graph-specific
obstruction has a name and an algebra. Neighbours fire at different times, so the second layer
integrates a locally reordered pre-activation; writing that down yields an exact error identity, and
the identity yields four things in turn: a closed-form forward bias $\betastar = \theta/(2T)$ that
removes the systematic part of the error and makes the method parameter-free; a membrane readout
that removes the output quantisation and clipping that counting imposes, and that we prove---and
verify to four decimals---is exactly invariant to the order of upstream spikes, which the count
decoder is not; a margin-recovery theorem whose required latency depends on neighbourhood geometry
rather than on graph size; and, once the identity is extended to $L$ layers, a depth law: the error
is amplified by the cumulative spectral gain $G=\prod_l\lVert M_l\rVert_2$ of the trained ANN, and
the amplification follows $G^{0.95}$ with $r=0.990$ over nine cells, so what a deep model will cost
in latency is computable from its weights before any simulation.

At seven seeds across six graphs, converted networks reach their ANN to within $1.5$ points at
$T \le 8$, and on a protocol shared by four arms they are never dominated: at matched accuracy they
cost $2$--$23\times$ fewer synaptic operations than a trained spiking network on two of three graphs
under the static-update accounting --- the ordering is convention-conditional
(\secref{sec:accounting}), which is itself one of the paper's findings --- and on CiteSeer they are
the only arm that reaches the band at all. The identity also buys a concrete decision: in the deep
regime, where uniform budgets explode with the depth law, a validation-only per-layer window search
keeps the accuracy band with $22.7$--$93.4\%$ fewer layer-updates on six of eight held-out
four-layer cells (\secref{sec:allocation}), while the weight-only routes fail to transfer
(\secref{sec:transfer}). External anchors frame, and do not enter, this protocol: SpikeGraphormer's
released pipeline reaches $74.0$ on Cora below the conversion at $T=16$, and SpikeGCL's
self-supervised pipeline reaches $86.6$ above every arm here. The converted networks train one
model instead of five, and they degrade in the way the theory predicts when the aggregation
operator, the threshold, the width or the decoder is changed.

The limits are as informative as the results. Conversion does not
buy energy: at matched accuracy it emits more spikes than a trained spiking network, and pruning
half the hidden units recovers only a third of the excess at a fraction of a point of accuracy.
Depth is not bought either: the amplification cannot be reparametrised away, because the spike
threshold is an absolute scale while the ReLU that would let one rescale it is homogeneous---and at
$L\ge3$ the closed form that makes the two-layer case free of simulation breaks too, for the same
reason, the reset. The count decoder does not free the computation from temporal ordering, since
only the membrane decoder has that. And averaging over the threshold and the initial
potential---axes that cost no spikes---never beats the ANN, whereas nine independently trained ANNs
beat it by about a point on every graph, so what helps is ensembling models and not re-sampling one
model's discretisation error. The published numbers cannot be used as baselines without checking:
the two anchors we could place beside our own, measured from their authors' code on the public
split, bracket the conversion rather than rank it.

A prediction that failed is reported alongside them. The laws rank candidate models almost
perfectly ($\rho=0.967$) but miss the absolute conversion gap by $4.6$ points, and the ablation that
would have confirmed the mechanism does not. What survives is narrower and more useful
than a transferable predictor: conversion is a statement about how a computation is discretised,
on a graph the discretisation happens along the neighbourhood as well as along time, and the two
quantities that decide whether a given model converts well---its activation mass and its cumulative
spectral gain---can both be read off the trained weights before a single spike is simulated.
